\chapter{Implementation}

\section{Project Scope and Evolution}

The project was initially scoped as a dual-track investigation: a \textbf{self-trained track}, in which an OSNet model would be trained from scratch on the Market-1501 training set and then evaluated on anonymized variants, and a \textbf{pretrained track}, using established pretrained weights from \texttt{torchreid}. The self-trained track was intended to investigate whether a model trained on original data would exhibit different sensitivity to anonymization compared to a model trained with standardized augmentation and hyperparameter tuning.

During development, the self-trained track deviated significantly from the core research question. Training a \acrshort{reid} model from scratch introduced a substantial set of secondary concerns: hyperparameter tuning (learning rate schedules, batch sizes, sampling strategies), architecture selection, training data augmentation, convergence validation, and the need to establish that the self-trained model was a fair comparator before any anonymization analysis could begin. These issues consumed development time without contributing directly to the privacy--utility trade-off investigation, which is the primary focus of this project.

The decision was made to \textbf{rescope to the pretrained track only}. This refocusing provided several advantages: (a)~the pretrained \texttt{osnet\_x1\_0} weights from \texttt{torchreid} represent a well-validated, reproducible baseline that isolates the anonymization effect from model training noise, (b)~development effort could be redirected toward building a comprehensive four-level anonymization pipeline and a thorough evaluation framework, and (c)~the results are directly comparable to published \acrshort{reid} benchmarks that use the same pretrained model.

Table~\ref{tab:timeline} summarizes the project phases.

\begin{table}[H]
    \centering
    \begin{tabular}{@{}clp{8cm}@{}}
        \toprule
        \rowcolor{cfTeal} \thdr{Phase} & \thdr{Title} & \thdr{Description} \\
        \midrule
        1 & Initial Scoping & Defined research questions. Planned dual-track approach (self-trained + pretrained OSNet). Identified Market-1501 as the benchmark dataset. \\
        2 & Anonymization Pipeline & Implemented L1 (blur) and L4 (edge) as CPU-bound methods. Developed the resolution-handling pipeline. Implemented L2 (SD Inpainting) and L3 (RAD) on Kaggle GPU notebooks. \\
        3 & Scope Revision & Self-trained track encountered significant scope deviation. Rescoped to pretrained OSNet only to focus on the core privacy--utility trade-off analysis. \\
        4 & Evaluation Framework & Deployed pretrained OSNet for utility evaluation. Trained MobileNetV2 attacker for privacy evaluation. Regenerated L1 dataset with upgraded MediaPipe face detection. \\
        5 & Integration & Built Streamlit dashboard. Compiled trade-off analysis, visualizations, and final report. \\
        \bottomrule
    \end{tabular}
    \caption{Project development phases.}
    \label{tab:timeline}
\end{table}


\section{Development Environment}

All experiments were conducted on a system with an AMD Ryzen 5 7600X CPU (6 cores), AMD Radeon RX 7800 XT GPU (16~GB VRAM), Python 3.12, and PyTorch 2.11.0 with ROCm/HIP backend (\texttt{device="cuda"} mapped transparently to HIP via the ROCm compatibility layer). GPU-bound dataset generation (Levels~2 and~3) was initially performed on Kaggle's free GPU tier (NVIDIA Tesla P100, 16~GB VRAM) before local hardware was available. Table~\ref{tab:dependencies} lists the principal libraries.

\begin{table}[H]
    \centering
    \begin{tabular}{@{}lll@{}}
        \toprule
        \rowcolor{cfTeal} \thdr{Library} & \thdr{Version} & \thdr{Purpose} \\
        \midrule
        torchreid & 0.2.5 & ReID model management and evaluation \\
        diffusers & 0.37.1 & Stable Diffusion pipelines \\
        controlnet-aux & 0.0.10 & OpenPose skeleton extraction \\
        mediapipe & 0.10.35 & Face detection (Level~1) \\
        torchvision & 0.26.0 & MobileNetV2 and image transforms \\
        opencv-python & 4.13.0 & Image processing and Canny edge detection \\
        streamlit & 1.45.0 & Interactive dashboard \\
        plotly & 6.1.2 & Interactive chart visualizations \\
        \bottomrule
    \end{tabular}
    \caption{Principal software dependencies.}
    \label{tab:dependencies}
\end{table}


\section{Dataset Generation}

The four anonymized dataset variants were generated by applying each anonymization function to all images in the three Market-1501 splits (\texttt{bounding\_box\_train}, \texttt{bounding\_box\_test}, and \texttt{query}), producing approximately 36,000 anonymized images per level. Each anonymized dataset preserves the original Market-1501 directory structure and filename conventions, ensuring compatibility with the \texttt{torchreid} data loading pipeline.

Gaussian blur (Level~1) and Canny edge detection (Level~4) are CPU-bound operations and were initially generated on Kaggle's free tier in approximately 15 minutes each. The Level~1 dataset was subsequently regenerated locally using the upgraded MediaPipe face detection implementation, which replaced an earlier OpenCV Haar cascade approach that exhibited lower detection accuracy on the small 64$\times$128 images.

Stable Diffusion inpainting (Level~2) and RAD (Level~3) require \acrshort{gpu} acceleration and were generated using Kaggle \acrshort{gpu} notebooks. These levels apply the upscale--process--downscale pipeline described in Section~\ref{sec:resolution}, processing each image individually through the generative model. Due to the sequential nature of diffusion inference, generation of each GPU-bound dataset required several hours of processing time across the full 36,036 images.


\section{Attacker Model Training}

The MobileNetV2 identity classifier was fine-tuned on the original (non-anonymized) \texttt{bounding\_box\_test} images to learn the visual appearance of each of the 750 test/query identities. Table~\ref{tab:attacker_config} details the training configuration.

\begin{table}[H]
    \centering
    \begin{tabular}{@{}ll@{}}
        \toprule
        \rowcolor{cfTeal} \thdr{Parameter} & \thdr{Value} \\
        \midrule
        Base model & MobileNetV2 (ImageNet pre-trained) \\
        Training data & \texttt{bounding\_box\_test} (19,732 images, 750 IDs) \\
        Input resolution & 256$\times$128 pixels \\
        Final layer & Linear(1280, 750) \\
        Loss function & CrossEntropyLoss \\
        Optimizer & Adam ($\text{lr} = 3 \times 10^{-4}$) \\
        Batch size & 64 \\
        Max epochs & 30 \\
        Early stopping & Patience = 4 (on training loss) \\
        \bottomrule
    \end{tabular}
    \caption{MobileNetV2 attacker training configuration.}
    \label{tab:attacker_config}
\end{table}

Identity labels are parsed from Market-1501 filenames (the first four characters encode the person ID). Images with IDs $-1$ (junk) and $0000$ (distractor) are excluded from training. The final layer of MobileNetV2 is replaced with a linear classifier mapping the 1280-dimensional feature vector to 750 identity classes. All layers are fine-tuned (not frozen), allowing the ImageNet-pretrained features to adapt to the pedestrian identity classification task.


\section{Interactive Dashboard}

A Streamlit-based dashboard (\texttt{app.py}) provides interactive exploration of the experimental results across three tabs. The \textbf{Trade-off Analysis} tab presents dual-axis charts comparing mAP and privacy leakage, scatter plots mapping the privacy--utility space, and per-level metric summary cards. The \textbf{Visual Comparison} tab displays side-by-side query images across all anonymization levels, allowing visual inspection of how each technique transforms the original image. The \textbf{Detailed Metrics} tab provides \acrshort{cmc} curve charts, mAP waterfall diagrams, a heatmap of all metrics across levels, and a full data table for export. All charts use Plotly with the \texttt{plotly\_white} template and the Cassette Futurism color palette for visual consistency with this report.
